% GENERATED FILE. Edit canonical lessons, then run npm run generate:books.
% canonical-source-hash: fnv1a64:eb4f9f091bc93ed4
% canonical-lessons: KA-C537-kavi, KA-C537-kavite, KA-C537-sahitya, KA-C537-kalakrti, KA-C537-racisu

\chapter{Poetry and Art}
\label{ch:ka-poetry-and-art}
\hlchaptermodality{537}

\begin{chapteropening}
\textbf{By the end of this chapter:} \emph{I can say a poet, a poem, literature, a work of art and to compose.}

Complete the last lesson of chapter 537: I can say a poet, a poem, literature, a work of art and to compose.
\end{chapteropening}

\section[\texorpdfstring{\kn{ಕವಿ} (kavi)}{ಕವಿ (kavi)}]{\kn{ಕವಿ} (kavi) — a poet}

\label{lesson:KA-C537-kavi}



\begin{quote}
Before the new one: say the Kannada for to perform, then the Kannada for to act.
\end{quote}

\subsection*{You'll want to know: \kn{ಕವಿ}}
\textbf{\kn{ಕವಿ}} — \emph{kavi} — \textquotedblleft{}a poet\textquotedblright{}.

\begin{quote}
\kn{ಆ} \kn{ಕವಿ} \kn{ಮಳೆಯ} \kn{ಬಗ್ಗೆ} \kn{ಸುಂದರ} \kn{ಹಾಡು} \kn{ಬರೆದರು}.
\end{quote}

\begin{quote}
\emph{ā kavi maḷeya bagge sundara hāḍu baredaru.}
\end{quote}

\begin{quote}
That poet wrote a beautiful song about the rain.
\end{quote}

Plural: \kn{ಕವಿಗಳು}, \emph{kavigaḷu}.

A woman poet is \kn{ಕವಯಿತ್ರಿ}.

\begin{cousinweb}
\textbf{\kn{ಕವಿ}} (\emph{kavi}) is borrowed from Sanskrit \emph{kavi}, \textquotedblleft{}wise; a wise man, a poet\textquotedblright{} (CDIAL 2964).
\end{cousinweb}

\begin{grammarlens}[title={What you've built}]
One more word for describing what you lived and felt.
\end{grammarlens}

\subsection*{Your turn}
\begin{itemize}
\raggedright
  \item \emph{Say it:} \emph{kavi}
  \item \emph{Say it:} \emph{kavi}, once more
  \item \emph{Say it:} \emph{ā kavi maḷeya bagge sundara hāḍu baredaru.} — That poet wrote a beautiful song about the rain.
  \item \emph{Recall:} say the Kannada for to perform, then the Kannada for to act, then say \emph{kavi} again
\end{itemize}

\begin{tcolorbox}[breakable,colback=teal!4,colframe=teal!35!black,title={Before you move on}]
What does \kn{ಕವಿ} mean? (\textquotedblleft{}A poet\textquotedblright{}.) Say it once more.
\end{tcolorbox}
\section[\texorpdfstring{\kn{ಕವಿತೆ} (kavite)}{ಕವಿತೆ (kavite)}]{\kn{ಕವಿತೆ} (kavite) — a poem}

\label{lesson:KA-C537-kavite}



\begin{quote}
Before the new one: say the Kannada for to act, then the Kannada for a poet.
\end{quote}

\subsection*{You'll want to know: \kn{ಕವಿತೆ}}
\textbf{\kn{ಕವಿತೆ}} — \emph{kavite} — \textquotedblleft{}a poem\textquotedblright{}.

\begin{quote}
\kn{ಮಗಳು} \kn{ಶಾಲೆಯಲ್ಲಿ} \kn{ಒಂದು} \kn{ಕವಿತೆ} \kn{ಓದಿದಳು}.
\end{quote}

\begin{quote}
\emph{magaḷu śāleyalli ondu kavite ōdidaḷu.}
\end{quote}

\begin{quote}
My daughter read a poem at school.
\end{quote}

Also \kn{ಕವನ}.

\begin{cousinweb}
\textbf{\kn{ಕವಿತೆ}} (\emph{kavite}) is borrowed from Sanskrit \emph{kavitā}, \textquotedblleft{}poetry\textquotedblright{} (Monier-Williams), built on \emph{kavi}, \textquotedblleft{}a poet\textquotedblright{}.
\end{cousinweb}

\begin{grammarlens}[title={What you've built}]
One more word for describing what you lived and felt.
\end{grammarlens}

\subsection*{Your turn}
\begin{itemize}
\raggedright
  \item \emph{Say it:} \emph{kavite}
  \item \emph{Say it:} \emph{kavite}, once more
  \item \emph{Say it:} \emph{magaḷu śāleyalli ondu kavite ōdidaḷu.} — My daughter read a poem at school.
  \item \emph{Recall:} say the Kannada for to act, then the Kannada for a poet, then say \emph{kavite} again
\end{itemize}

\begin{tcolorbox}[breakable,colback=teal!4,colframe=teal!35!black,title={Before you move on}]
What does \kn{ಕವಿತೆ} mean? (\textquotedblleft{}A poem\textquotedblright{}.) Say it once more.
\end{tcolorbox}
\section[\texorpdfstring{\kn{ಸಾಹಿತ್ಯ} (sāhitya)}{ಸಾಹಿತ್ಯ (sāhitya)}]{\kn{ಸಾಹಿತ್ಯ} (sāhitya) — literature; lyrics}

\label{lesson:KA-C537-sahitya}



\begin{quote}
Before the new one: say the Kannada for a poet, then the Kannada for a poem.
\end{quote}

\subsection*{You'll want to know: \kn{ಸಾಹಿತ್ಯ}}
\textbf{\kn{ಸಾಹಿತ್ಯ}} — \emph{sāhitya} — \textquotedblleft{}literature; lyrics\textquotedblright{}.

\begin{quote}
\kn{ಕನ್ನಡ} \kn{ಸಾಹಿತ್ಯ} \kn{ತುಂಬಾ} \kn{ಹಳೆಯದು}.
\end{quote}

\begin{quote}
\emph{kannaḍa sāhitya tumbā haḷeyadu.}
\end{quote}

\begin{quote}
Kannada literature is very old.
\end{quote}

Also the words of a song.

\begin{cousinweb}
\textbf{\kn{ಸಾಹಿತ್ಯ}} (\emph{sāhitya}) is borrowed from Sanskrit \emph{sāhitya}, \textquotedblleft{}association, union\textquotedblright{}, from \emph{sahita}, \textquotedblleft{}joined\textquotedblright{} (Monier-Williams): words and meaning joined well.
\end{cousinweb}

\begin{grammarlens}[title={What you've built}]
One more word for describing what you lived and felt.
\end{grammarlens}

\subsection*{Your turn}
\begin{itemize}
\raggedright
  \item \emph{Say it:} \emph{sāhitya}
  \item \emph{Say it:} \emph{sāhitya}, once more
  \item \emph{Say it:} \emph{kannaḍa sāhitya tumbā haḷeyadu.} — Kannada literature is very old.
  \item \emph{Recall:} say the Kannada for a poet, then the Kannada for a poem, then say \emph{sāhitya} again
\end{itemize}

\begin{tcolorbox}[breakable,colback=teal!4,colframe=teal!35!black,title={Before you move on}]
What does \kn{ಸಾಹಿತ್ಯ} mean? (\textquotedblleft{}Literature; lyrics\textquotedblright{}.) Say it once more.
\end{tcolorbox}
\section[\texorpdfstring{\kn{ಕಲಾಕೃತಿ} (kalākṛti)}{ಕಲಾಕೃತಿ (kalākṛti)}]{\kn{ಕಲಾಕೃತಿ} (kalākṛti) — a work of art}

\label{lesson:KA-C537-kalakrti}



\begin{quote}
Before the new one: say the Kannada for a poem, then the Kannada for literature.
\end{quote}

\subsection*{You'll want to know: \kn{ಕಲಾಕೃತಿ}}
\textbf{\kn{ಕಲಾಕೃತಿ}} — \emph{kalākṛti} — \textquotedblleft{}a work of art\textquotedblright{}.

\begin{quote}
\kn{ವಸ್ತುಸಂಗ್ರಹಾಲಯದಲ್ಲಿ} \kn{ಹಳೆಯ} \kn{ಕಲಾಕೃತಿಗಳು} \kn{ಇವೆ}.
\end{quote}

\begin{quote}
\emph{vastusaṅgrahālayadalli haḷeya kalākṛtigaḷu ive.}
\end{quote}

\begin{quote}
There are old works of art in the museum.
\end{quote}

Plural: \kn{ಕಲಾಕೃತಿಗಳು}, \emph{kalākṛtigaḷu}.

\kn{ಕಲೆ}, art + \kn{ಕೃತಿ}, a work.

\begin{cousinweb}
\textbf{\kn{ಕಲಾಕೃತಿ}} (\emph{kalākṛti}) is a learned compound built from Sanskrit \emph{kalā}, \textquotedblleft{}any practical art\textquotedblright{} (CDIAL 2927), and \emph{kṛti}, \textquotedblleft{}a work, a creation\textquotedblright{} (Monier-Williams).
\end{cousinweb}

\begin{grammarlens}[title={What you've built}]
One more word for describing what you lived and felt.
\end{grammarlens}

\subsection*{Your turn}
\begin{itemize}
\raggedright
  \item \emph{Say it:} \emph{kalākṛti}
  \item \emph{Say it:} \emph{kalākṛti}, once more
  \item \emph{Say it:} \emph{vastusaṅgrahālayadalli haḷeya kalākṛtigaḷu ive.} — There are old works of art in the museum.
  \item \emph{Recall:} say the Kannada for a poem, then the Kannada for literature, then say \emph{kalākṛti} again
\end{itemize}

\begin{tcolorbox}[breakable,colback=teal!4,colframe=teal!35!black,title={Before you move on}]
What does \kn{ಕಲಾಕೃತಿ} mean? (\textquotedblleft{}A work of art\textquotedblright{}.) Say it once more.
\end{tcolorbox}
\section[\texorpdfstring{\kn{ರಚಿಸು} (racisu)}{ರಚಿಸು (racisu)}]{\kn{ರಚಿಸು} (racisu) — to compose; to create}

\label{lesson:KA-C537-racisu}



\begin{quote}
Before the new one: say the Kannada for literature, then the Kannada for a work of art.
\end{quote}

\subsection*{You'll want to know: \kn{ರಚಿಸು}}
\textbf{\kn{ರಚಿಸು}} — \emph{racisu} — \textquotedblleft{}to compose; to create\textquotedblright{}.

\begin{quote}
\kn{ಈ} \kn{ಹಾಡನ್ನು} \kn{ಒಬ್ಬ} \kn{ಹಳ್ಳಿಯ} \kn{ಕವಿ} \kn{ರಚಿಸಿದರು}.
\end{quote}

\begin{quote}
\emph{ī hāḍannu obba haḷḷiya kavi racisidaru.}
\end{quote}

\begin{quote}
A village poet composed this song.
\end{quote}

Past (I …): \kn{ರಚಿಸಿದೆ}, \emph{raciside}.

\kn{ರಚನೆ}, a composition.

\begin{cousinweb}
\textbf{\kn{ರಚಿಸು}} (\emph{racisu}) is built on the Sanskrit root \emph{rac}, \textquotedblleft{}to arrange, to compose\textquotedblright{}, seen in \emph{racanā} (Monier-Williams), with the verb ending -\emph{isu}.
\end{cousinweb}

\begin{grammarlens}[title={What you've built}]
One more word for describing what you lived and felt.
\end{grammarlens}

\subsection*{Your turn}
\begin{itemize}
\raggedright
  \item \emph{Say it:} \emph{racisu}
  \item \emph{Say it:} \emph{racisu}, once more
  \item \emph{Say it:} \emph{ī hāḍannu obba haḷḷiya kavi racisidaru.} — A village poet composed this song.
  \item \emph{Recall:} say the Kannada for literature, then the Kannada for a work of art, then say \emph{racisu} again
\end{itemize}

\begin{tcolorbox}[breakable,colback=teal!4,colframe=teal!35!black,title={Before you move on}]
What does \kn{ರಚಿಸು} mean? (\textquotedblleft{}To compose; to create\textquotedblright{}.) Say it once more.
\end{tcolorbox}
